
Iterative code refinement with feedback improves code generation quality, but controlled comparisons between execution feedback and AI critique remain limited due to method-level confounds in prior work. This study presents a controlled comparison isolating feedback signals: using identical base models, refinement protocols, and benchmarks while varying only the feedback source and granularity. Experiments on 1,095 buggy code samples reveal that 84.7\% of execution traces contain counterfactual localization information (CF-score $\geq$ 0.4), with type errors (0.665) and logic errors (0.583) providing the highest information density. Targeted edits ($\leq 5$ lines changed) achieve a 68.4\% fix rate compared to 31.2\% for global rewrites (>5 lines). On a minimal validation set (n=5), detailed execution feedback achieves 100\% pass rate after refinement while binary pass/fail feedback achieves 60\%. Contrary to initial predictions, execution feedback advantage does not increase with task complexity. These findings suggest that the relevant distinction for feedback effectiveness is information structure---specifically, localization information---rather than information source.
